% M10 SSH-solver campaign — interim report with figures.
% Build:  module load texlive/live2025-gcc-13.3.0 && make
\documentclass[11pt,a4paper]{article}

\usepackage[utf8]{inputenc}
\usepackage[T1]{fontenc}
\usepackage{lmodern}
\usepackage[margin=22mm,top=18mm,bottom=18mm]{geometry}
\usepackage{amsmath,amssymb}
\usepackage{booktabs}
\usepackage{graphicx}
\usepackage{microtype}
\usepackage{parskip}
\usepackage{multirow}
\usepackage{titlesec}
\usepackage[hidelinks]{hyperref}

\titleformat{\section}{\normalfont\large\bfseries}{\thesection.}{0.6em}{}
\titlespacing*{\section}{0pt}{13pt}{5pt}
\titleformat{\subsection}{\normalfont\bfseries}{\thesubsection.}{0.5em}{}
\titlespacing*{\subsection}{0pt}{9pt}{3pt}
\setlength{\tabcolsep}{5pt}

\title{\vspace{-12mm}\large\bfseries Communication-avoiding solvers for the sea-surface-height
system in FESOM2/Kokkos}
\author{\normalsize Interim report}
\date{\normalsize 6 August 2026 --- Levante (128-core CPU nodes, A100 GPU nodes)}

\begin{document}
\maketitle
\vspace{-6mm}

\section{Scope}

The barotropic sea-surface-height (SSH) solve is the model's synchronisation bottleneck.
Preconditioned conjugate gradients perform two global reductions and two halo exchanges per
iteration, and their cost grows as a mesh is spread over more cores. We implemented four
alternatives behind a runtime switch --- Chronopoulos--Gear CG (\texttt{cg2}), pipelined CG
(\texttt{pipecg}), PIPECG-OATI (\texttt{oati}) and Chebyshev P-CSI (\texttt{pcsi}) --- and
measured them on four meshes and both backends. The certified baseline remains bit-identical
when the switch is unset, which two byte-level gates confirm.

Every timing reported here comes from a same-allocation A/B against the baseline solver,
minimum of two repetitions, 300 steps, with all four legs verified to have fired no fallback.
A leg that falls back is a mixture of the variant and the baseline and is not an A/B point;
this is the criterion under which an earlier set of fArc results was withdrawn.

\section{Two defects had to be repaired first}

\subsection{The preconditioner is not symmetric}

The preconditioner inherited from FESOM's \texttt{solver.F90} is $D^{-1}C$ with $C$ symmetric.
The leading $1/d_i$ breaks symmetry, with a measured defect ratio of 0.62--0.74 across three
meshes. Every CG--CG variant replaces the explicit inner product $(p,Ap)$ with a recurrence
that is valid only for a symmetric preconditioner, so the step length came out wrong by
21.8\,\% and most solves diverged. Symmetrising to $D^{-1/2}CD^{-1/2}$ removes the error.

The same repair shortens the \emph{existing} solver by about 7\,\% in iterations (126 to 117
on CORE2), so correcting \texttt{solver.F90} would benefit the Fortran model independently of
anything else here.

\subsection{The fallback guard was aborting healthy solves}

Our own guard declares failure after 20 consecutive iterations without a 0.1\,\% residual
drop. On fArc it fired at iteration 107 of solves that the baseline completes in 212, and it
was the reason every fArc result was withdrawn. Re-running the identical dumped system with
the guard widened, \texttt{cg2} converges at iteration 205 --- six \emph{fewer} than the
baseline's 211. The residual falls again three iterations after the point where the guard had
given up.

The threshold is below what the reference method itself requires: baseline CG's own longest
plateau on fArc is 21 iterations. Residual non-monotonicity is ordinary CG behaviour, since
only the $A$-norm of the error decreases monotonically, and on a system with condition number
843 a 20-iteration plateau is unremarkable. Baseline CG never reported a fallback for a
simpler reason than robustness: it carries no guard at all, so its counter cannot increment.

\section{Results}

\begin{figure}[t]
\centering
\includegraphics[width=\textwidth]{/work/ab0995/a270088/port2/m10/figs/fig_m10_payoff.png}
\caption{\small \textbf{(a)} Whole-step speed-up against the share of the step spent in the
SSH solve, one point per configuration, best solver at that configuration. Filled circles are
CPU, open triangles GPU. \textbf{(b)} The speed-up of the solve phase itself for \texttt{oati};
hatched bars are GPU. The two backends separate in (b), which is what makes them separate in
(a): on CORE2 and fArc the solve is 46--47\,\% faster on CPU but only 18--21\,\% faster on GPU.
On DARS, where the solve is a small share of the step, it is 27\,\% faster on CPU and 12\,\%
\emph{slower} on GPU.}
\label{fig:payoff}
\end{figure}

\subsection{The payoff tracks the SSH share of the step}

Figure~\ref{fig:payoff}(a) collects every valid configuration. The whole-step gain follows the
share of the step the solve occupies, which varies from 0.8\,\% (DARS at 512 CPU ranks) to
42.2\,\% (fArc on 16 GPU nodes). Below roughly 5\,\% the effect is within measurement noise;
above it the gain grows steadily.

That share is not a property of the mesh but of how thinly it is spread. On CORE2 it rises
from 2.5\,\% at 128 ranks to 18.9\,\% at 864. The solve is also the only part of the model
that anti-scales: from 128 to 864 ranks everything else speeds up $5.40\times$ (79\,\% parallel
efficiency) while the solve becomes 41\,\% \emph{slower} in absolute terms. The solvers
therefore pay most exactly where the model stops scaling, and on CORE2 they move the
parallel-efficiency knee out by a factor of 1.26--1.29.

\begin{center}\small
\begin{tabular}{llrrrr}
\toprule
& configuration & vertices/rank & SSH\,\% & best solver & whole-step \\
\midrule
\multirow{4}{*}{CPU}
& CORE2, 512 ranks (production) & 248 & 9.9 & \texttt{pcsi} & $-6.30$\,\% \\
& CORE2, 864 ranks & 146 & 18.9 & \texttt{pcsi} & $-13.10$\,\% \\
& fArc, 2048 ranks & 312 & 22.8 & \texttt{oati} & $-13.28$\,\% \\
& DARS, 8192 ranks & 385 & 6.4 & \texttt{pcsi} & $-4.04$\,\% \\
\midrule
\multirow{4}{*}{GPU}
& CORE2, 2 nodes & 15857 & 25.8 & \texttt{oati} & $-5.32$\,\% \\
& CORE2, 16 nodes & 1982 & 30.5 & \texttt{oati} & $-6.73$\,\% \\
& fArc, 16 nodes & 9975 & 42.2 & \texttt{oati} & $-10.36$\,\% \\
& DARS, 16 nodes & 49380 & 7.4 & \texttt{oati} & $-1.00$\,\% \\
\bottomrule
\end{tabular}
\end{center}

\subsection{Which solver}

\texttt{oati} wins at every configuration measured in range on both backends and tracks the
baseline's iteration count everywhere, which makes it the safer default. \texttt{pcsi}
converges faster where its Chebyshev interval is accurate --- it is the best solver on CORE2
CPU --- but that interval rests on an eigenvalue estimate that was poor on fArc, where
\texttt{pcsi} costs 8--10\,\% on both backends while needing 376 iterations against 212.
\texttt{pipecg} is indistinguishable from \texttt{cg2} on this machine, because no
\texttt{MPI\_Iallreduce} progression is available on any Levante stack we tested.

P-CSI performs no reduction per iteration at all, its only global communication being the
convergence test every $K$ iterations. Raising $K$ from 1 to the shipped default of 5 more than
doubles the whole-step gain, from $-4.86$\,\% to $-11.11$\,\% at CORE2 864 ranks. Beyond 5 the
return flattens and reappears as wasted work: $K=20$ buys 1.2 percentage points more while
overshooting the tolerance by 6.7\,\% more iterations.

\subsection{How much of this depends on the initial condition}

Because the fill defect above changes every cold-start trajectory, nine of the configurations
were re-measured with the deterministic fill in place. The comparison is between ratios, each
of which is internally consistent --- all four legs of a comparison share one allocation and
one initial condition --- so it asks whether the conclusions survive a different, correct
starting state rather than whether the timings repeat.

\begin{center}\small
\begin{tabular}{llrrr}
\toprule
configuration & solver & legacy fill & deterministic fill & iterations \\
\midrule
fArc, 2048 CPU ranks  & \texttt{oati} & $-13.28$\,\% & $-12.90$\,\% & $211 \to 213$ \\
fArc, 1024 CPU ranks  & \texttt{oati} & $-5.60$\,\%  & $-5.29$\,\%  & $212 \to 213$ \\
fArc, 16 GPU nodes    & \texttt{oati} & $-10.36$\,\% & $-10.58$\,\% & $212 \to 213$ \\
DARS, 8192 CPU ranks  & \texttt{oati} & $-3.94$\,\%  & $-3.32$\,\%  & $36.7 \to 36.7$ \\
DARS, 6144 CPU ranks  & \texttt{oati} & $-3.01$\,\%  & $-1.69$\,\%  & $36.6 \to 36.6$ \\
CORE2, 1 GPU node     & \texttt{pcsi} & $-5.19$\,\%  & $-3.11$\,\%  & $121 \to 133$ \\
NG5, 16 GPU nodes     & \texttt{pcsi} & $-2.32$\,\%  & $-1.70$\,\%  & $78 \to 87$ \\
\bottomrule
\end{tabular}
\end{center}

\noindent
The headline results stand: fArc reproduces on both backends, including the campaign's best
figure of $-10.4$ against $-10.6$\,\% at 16 GPU nodes. Two effects are worth stating plainly.

On DARS the gain shrinks, and the reason is visible in the per-iteration cost. The variants are
unchanged in absolute terms, while the baseline solve becomes 14--19\,\% cheaper at an unchanged
iteration count --- that is, it waits less. Since roughly nine tenths of the measured solve cost
at these scales is time spent waiting at collectives, and the baseline is the most synchronous
of the four methods, a clean initial condition removes imbalance that was penalising it
specifically. The low-synchronisation solvers did not get worse; the method they are measured
against got better.

The second effect is confined to P-CSI and is a caution about how its numbers travel. Chebyshev
iteration counts are fixed by a spectral estimate made at setup from the operator, which depends
on the state; where that estimate moved, so did the result --- 9.5\,\% more iterations at CORE2
GPU and 11.5\,\% at NG5 GPU, enough to cost P-CSI first place at the former and most of its lead
at the latter --- while at fArc and DARS, where the estimate did not move, its numbers are
unchanged. This is the same sensitivity that made a deliberately truncated Lanczos run look
fast earlier in the campaign. It does not weaken P-CSI as an option, but a P-CSI figure should
be quoted together with the state its interval was estimated on, in a way a \texttt{cg2} or
\texttt{oati} figure need not be.

\section{Where the step actually goes}

\begin{figure}[t]
\centering
\includegraphics[width=\textwidth]{/work/ab0995/a270088/port2/m10/figs/fig_m10_budget.png}
\caption{\small Per-phase compute and MPI wait, baseline against the best solver, from
per-rank wall clocks with wait captured inside MPI by PMPI interposition. About half of each
step is wait. The SSH bar is almost entirely wait, and the variant also shortens the wait in
forcing and ice dynamics, phases whose own MPI call counts are unchanged.}
\label{fig:budget}
\end{figure}

Figure~\ref{fig:budget} resolves the step into compute and MPI wait. Wait accounts for 51\,\%
of the step on CORE2 at 864 ranks and 53\,\% on fArc at 2048: the model is
synchronisation-bound at these rank counts. The SSH solve is the single largest source of
wait, but roughly 90\,\% of the solve's own cost \emph{is} wait --- its arithmetic is 1.0\,ms
on CORE2 and 2.5\,ms on fArc.

Compute is unchanged between baseline and variant (23.5 to 23.1\,ms on CORE2), so the variants
buy nothing arithmetically and the entire gain is synchronisation. Only 5.8 of the 9.5\,ms
saved lies inside the solver: 3.7\,ms, or 39\,\%, appears in phases whose MPI call counts and
compute times are identical. Removing synchronisation points lets neighbouring phases stop
absorbing skew, which is why the whole-step gain exceeds the solve's own share of the step
rather than merely matching it.

\section{Load imbalance}

\begin{figure}[t]
\centering
\includegraphics[width=\textwidth]{/work/ab0995/a270088/port2/m10/figs/fig_m10_balance.png}
\caption{\small \textbf{(a)} Ocean-phase compute against the number of 3-D nodes each rank
owns, all 2048 fArc ranks. \textbf{(b)} Step time with a vertically balanced partition
relative to a surface-balanced one, at the same rank count in the same allocation.}
\label{fig:balance}
\end{figure}

The ocean phase is the largest block of real computation, and its per-rank cost varies by a
factor of 4.8. A bulk-synchronous phase costs its slowest rank, so this is a direct tax:
5.2\,ms of a 48.2\,ms step on CORE2 and 15.9\,ms of 84.7\,ms on fArc.

Figure~\ref{fig:balance}(a) identifies the cause. Ocean compute correlates with the number of
3-D nodes a rank owns at $r=0.967$, and with the number of 2-D nodes at $r=0.003$. The
partition balances surface nodes to within 1\,\%, but the work is per water column, and 3-D
nodes span a factor of 9.4. The imbalance is bathymetric. Solver wait correlates with 3-D
nodes at $-0.771$: deep-column ranks wait least because they are the stragglers, so the ranks
everyone waits for are the deep ones, not the ice-covered ones.

Balancing the vertical distribution does not pay. Weighting the partitioner by 3-D as well as
2-D nodes removes the imbalance completely (spread $9.6\times \to 1.05\times$) but forces the
partitioner to cut about ninety times more edges. The resulting fragmentation costs more
compute than the balance recovers: on CPU it is worth $-4.6$\,\% at 495 vertices per core,
breaks even near 250 and costs $+8.7$\,\% at 146; on GPU it costs $+29.7$\,\%, since the device
is far more sensitive to memory-access locality. Note that halo node counts understate the
damage --- they grow 40\,\% where the edge cut grows ninetyfold --- so the cut, not the halo,
is the quantity to watch.

\subsection{A mesh's shipped partitions are not all built the same way}

The partitions distributed with a mesh can differ in quality, and a poor one is not merely
slow: it can destabilise the model. NG5 illustrates both halves. Its shipped decompositions at
1024 and 2048 pieces balance 3-D nodes to within 4\,\%, while those at 4096 and 8192 pieces
are skewed by a factor of 13.9 --- the small counts were built with dual 2-D+3-D weighting and
the large ones were not.

That split predicts the model's behaviour exactly. Holding mesh, timestep, cold start, binary
and rank count fixed and varying only the decomposition, the shipped partition at 4096 pieces
reaches the model's blow-up guard at step 175 with velocities still climbing, while a freshly
generated one completes all 300 steps at 0.6205\,s/step:

\begin{center}\small
\begin{tabular}{lrl}
\toprule
NG5, 4096 ranks, dt 180, cold & 3-D balance & 300 steps \\
\midrule
shipped partition & $13.89\times$ & blow-up at step 175 \\
freshly generated & $1.05\times$ & completed \\
\bottomrule
\end{tabular}
\end{center}

\noindent
We had previously read this blow-up as a stability limit of the timestep, following an earlier
campaign that measured the same death at step 200 in this port and step 203 in Fortran. That
evidence came from the same shipped partition, so both readings attributed to the
configuration what belongs to one decomposition; Fortran failing there as well shows only that
Fortran dislikes it too.

The mechanism has since been identified, and it is not the balance. The climatology
hole-filling routine that supplies initial temperature and salinity where the 1$^\circ$
climatology has no data --- the entire Sea of Marmara, and most straits and marginal seas ---
fills each dummy node from its already-valid neighbours in local numbering order, running each
rank to exhaustion before the next halo exchange. The result therefore depends on the
decomposition. On an unlucky partition it writes a density front of order
10\,kg\,m$^{-3}$ across a single element in a strait separating water masses that differ by
some 20\,PSU; the resulting pressure-gradient force is constant, drives a linear velocity ramp
from the first step, and reaches the CFL limit a few hundred steps later. Every scheme dies on
such a partition. Balance was a marker only in the accidental sense that the same partitions
were built by different recipes. A deterministic fill --- ring-filling followed by relaxation,
with neighbour sums accumulated in global-index order --- removes the dependence: initial
conditions become bitwise identical across decompositions, and all four previously fatal NG5
partitions complete 300 steps. The same defect exists upstream in Fortran, where cold starts on
unlucky partitions inherit the same fragility (restarts are immune, which is why it is rarely
seen).

The practical rule is unchanged in form but now has a cause behind it: a cold-start blow-up
that appears on one decomposition and not another is an initial-condition artifact, not a
property of the timestep, the viscosity scheme or the solver.

\section{Status and limitations}

CPU and GPU coverage is complete for CORE2, fArc and DARS. NG5 now has two CPU rows as well,
measured with the deterministic fill described above: at 4096 ranks the solve is 1.9\,\% of the
step and the best variant returns $-0.69$\,\%, at 8192 ranks 3.6\,\% and $-1.93$\,\%. Both are
honest negatives and both sit where the share law puts them; NG5 is simply a mesh on which the
barotropic solve is a small part of the step at the rank counts we can reach. The in-range
partitions (16384, 20480 and 24576 pieces --- 451, 361 and 301 vertices per core) are generated
and usable, and the smallest of them needs 128 nodes; those rungs are left to the
split-explicit campaign, which is measuring the same ladder. NG5 on GPU rests on a single
16-node configuration.

All numbers come from one machine and a 300-step window. The default solver is unchanged: the
switch selects a variant only when set, and which variant to adopt --- together with the value
the widened guard threshold should ship with --- remains open.

\end{document}
